% =============================================================================
% Chapter 3: Scope
% =============================================================================

\chapter{Scope}
\label{ch:scope}

\section{In Scope}

\itc{} addresses the following problem domain:

\begin{itemize}
  \item \textbf{Multidimensional experimental data}: wind tunnel
    campaigns with named physical dimensions (angle of attack, sideslip,
    Mach number, Reynolds number, RPM, blade pitch, control-surface
    deflections, horizontal tail incidence, repeat number, etc.) and
    measured variables (forces, moments, pressures, temperatures).

  \item \textbf{CFD post-processing data} (surface $C_p$ distributions, propeller disk loads in polar coordinates, slipstream characterization) extracted from CFD solvers and organized along named dimensions.

  \item \textbf{Aerospace engineering computations}: performance analyses
    (takeoff, climb, envelope), flight dynamics (static margin, trim with
    uncertainty propagation, V-n diagram), and propulsion modeling
    (BEMT). Results are generated as VarFrame objects so the same
    data-management, plotting, comparison, and uncertainty machinery
    applies seamlessly to computed and measured data.

  \item \textbf{Multi-fidelity datasets}: combining wind tunnel and CFD
    data sources, with explicit fidelity tagging via dimensions and
    downstream surrogate-modeling support for cross-fidelity reconstruction.

  \item \textbf{Uncertainty quantification}: automatic propagation of
    measurement and modeling uncertainties through derived quantities via
    symbolic differentiation of the RPN expression tree (chain rule on the
    AST), following GUM JCGM~100:2008~\citep{gum2008}, including covariance
    terms when correlations are declared. Monte Carlo propagation is
    available specifically for models containing discrete branches (e.g.\
    conditional logic) where the linearization assumption breaks down.

  \item \textbf{Coordinate system handling}: canonical AIAA body axes by
    default, plus user-defined custom axes with rotation matrices
    (stability, wind, tunnel, propeller-disk), with rotation operations
    that propagate uncertainty correctly.

  \item \textbf{Provenance and History}: immutable records of where the
    data came from and what was done to it, with content hashing, user
    identification, and timestamps. Native export to PROV-N and PROV-JSON
    for interoperability with W3C-compliant data catalogs.

  \item \textbf{Reusable analysis workflows}: encapsulation of complete
    processing pipelines (corrections, derivations, uncertainty
    propagation) in named processors (\code{pproc}) and reusable equation
    files (\code{.itceq}).

  \item \textbf{Statistical analysis of repeated measurements}: collapse
    of a repeat dimension into mean, standard deviation, and 95\,\%
    confidence intervals.

  \item \textbf{Dataset comparison}: structured pairwise and N-way
    comparison of datasets with delta, percentage delta, and
    sigma-normalized differences using named inputs.

  \item \textbf{Publication-quality plotting}: AIAA-style figures
    generated declaratively, with semantic axis / color / symbol mapping,
    single-plot and multi-plot output, and origin-tag aware marker styling.

  \item \textbf{Surrogate modeling}: interface with the SMT
    library~\citep{smt} for Kriging and other surrogates, with VarFrame
    input/output and exportable standalone surrogate models.

  \item \textbf{Automated reporting}: LaTeX-backed PDF reports generated
    from processor results, with diagnostic plots, tabulated summaries,
    and warnings; the underlying \code{.tex} sources can be retained
    optionally.

  \item \textbf{Unit metadata and conversion}: optional unit annotations
    on variables and dimensions, with a custom \itc{}-internal conversion
    table; no external unit-library dependency.
\end{itemize}

\section{Out of Scope}

\begin{nreqbox}{NREQ-01}{General-purpose data science}
\itc{} is not a replacement for pandas, NumPy, or SciPy in general data
science. It is scoped to structured, named-dimension scientific datasets in
engineering contexts. Bridges to pandas (\code{to\_pandas}) and NumPy
(\code{to\_numpy}) are provided for interoperability.
\end{nreqbox}

\begin{nreqbox}{NREQ-02}{CFD solver or flow visualization}
\itc{} processes data extracted from CFD solvers. It does not interface with
solver native formats (e.g.\ OpenFOAM, SU2, Fluent), perform mesh
operations, or provide 3D flow visualization. It receives post-processed
scalar fields organized in named arrays.
\end{nreqbox}

\begin{nreqbox}{NREQ-03}{General-purpose plotting library}
The \itc{} plot system wraps matplotlib (default) and plotly (optional, for
interactive figures) for domain-specific convenience. It does not replace
either backend for general visualization tasks. Custom figures use the
\code{fig.mpl} escape hatch on the underlying \code{matplotlib.Figure}.
\end{nreqbox}

\begin{nreqbox}{NREQ-04}{Database or cloud storage backend}
\itc{} manages data as files on a local filesystem. It provides no database
connectivity, cloud storage integration, or remote dataset management.
The \code{.itc} binary format and CSV/JSON exports are the storage
mechanisms.
\end{nreqbox}

\begin{nreqbox}{NREQ-05}{Machine learning training pipelines}
The surrogate interface (\code{itc.surrogate}) is scoped to interpolation
and approximation of engineering datasets (Kriging, RBF, polynomial, IDW).
It is not a machine learning framework and does not support classification,
neural networks, or GPU-accelerated training.
\end{nreqbox}

\begin{nreqbox}{NREQ-06}{Real-time or streaming data acquisition}
\itc{} processes data from completed experiments stored in files. It does
not interface with DAQ hardware, real-time measurement streams, or live
sensor feeds.
\end{nreqbox}

\begin{nreqbox}{NREQ-07}{GUI or web application}
\itc{} is a Python library. It has no graphical user interface or web
dashboard. It is designed for use in Python scripts and Jupyter notebooks.
A browser-rendered diagnostic viewer is on the post-v0.1 roadmap, but it is
a read-only viewer, not an interactive application.
\end{nreqbox}

\begin{nreqbox}{NREQ-08}{Overloaded Python arithmetic on VarFrames}
\itc{} does not support \code{db1 + db2}, \code{db1 - db2}, \code{db1 * db2},
or \code{db1 / db2} via overloaded Python operators. All numerical
combinations of VarFrames go through \code{db.combine(other, op=...)}, which
records the operation in History and applies the declared correlation
structure. This is a deliberate decision (P-13).
\end{nreqbox}

\begin{nreqbox}{NREQ-09}{External unit-conversion libraries}
\itc{} does not depend on \code{pint}, \code{astropy.units}, or any other
external unit library. Unit metadata is optional and conversion is handled
by a custom internal table in \code{utils/units.py}.
\end{nreqbox}

\begin{nreqbox}{NREQ-10}{Solver execution and automation}
\itc{} does not launch, script, or automate any solver, panel code, or
simulation tool, and it embeds no solver-specific command emitters.
Dedicated driver packages own that responsibility and remain outside
\itc{}. Drivers may produce files or arrays that \code{itc.load} consumes,
and they may adopt \itc{} as their post-processing layer. This boundary
keeps \itc{} solver-agnostic and keeps every driver's version-compatibility
problem out of the \itc{} core.
\end{nreqbox}

\section{Positioning Among Existing Tools}
\label{sec:positioning}

A public engineering library must state where it sits relative to the
established scientific Python stack, both to justify its existence and to
make deliberate non-reuse auditable. Table~\ref{tab:positioning} maps the
nearest neighbors.

\begin{table}[H]
\centering
\caption{Positioning of \itc{} relative to existing tools.}
\label{tab:positioning}
\small
\begin{tabularx}{\textwidth}{@{}lXX@{}}
\toprule
\textbf{Tool} & \textbf{What it provides} & \textbf{Relation to \itc{}}\\
\midrule
xarray~\citep{xarray} &
Named dimensions and coordinates over NumPy arrays; the closest structural
relative of the VarFrame. &
\itc{} adopts the named-dimension data model but not the dependency: the
core stays NumPy-only (DD-02), and \itc{} adds mandatory provenance,
operating modes, origin tags, and GUM covariance propagation, which xarray
does not provide. Its alignment and broadcasting conventions inform the
\code{concat} and \code{combine} rules.\\
\addlinespace
pandas~\citep{pandas} &
Labeled tabular (2-D) data structures and IO. &
Different shape of problem: VarFrame is N-dimensional and grid-oriented.
A lazy bridge is provided in both directions (REQ-05, REQ-70).\\
\addlinespace
uncertainties~\citep{uncertaintiespkg} &
First-order uncertainty propagation with correlations through per-value
\code{ufloat} objects. &
Same GUM basis, different mechanics: per-value object wrapping conflicts
with the vectorized NumPy-only core and with the UncFrame mirror (DD-05).
\itc{} propagates over whole arrays symbolically on the expression tree.\\
\addlinespace
pint~\citep{pint} &
General unit parsing and conversion. &
Rejected as a dependency (DD-13, NREQ-09): the hand-curated internal table
keeps conversions auditable and avoids dtype interactions with NumPy and
SMT.\\
\addlinespace
SMT~\citep{smt} &
Surrogate modeling toolbox (Kriging, RBF, IDW). &
Adopted as the optional surrogate backend behind \code{itc.surrogate}
(REQ-63).\\
\bottomrule
\end{tabularx}
\end{table}

The differentiator of \itc{} is the combination, not any single feature:
mandatory provenance and history, operating modes with export guards,
origin tags, GUM propagation with covariance, and domain processors for
wind tunnel, CFD, and flight-test workflows. Several tools offer parts of
this; none offers the combination. Where an existing convention is good,
\itc{} mirrors it deliberately: named dimensions from xarray, docstring
standards from NumPy.

\section{Primary Use Cases}

Table~\ref{tab:use-cases} summarizes the primary use cases that drive the
requirements in this document. UC-01 through UC-06 cover experimental and
CFD post-processing scenarios; UC-07 through UC-09 cover aerospace
engineering computations and the multi-fidelity workflow.

\begin{table}[H]
\centering
\caption{Primary use cases driving \itc{} requirements.}
\label{tab:use-cases}
\small
\begin{tabularx}{\textwidth}{@{}llX@{}}
\toprule
\textbf{ID} & \textbf{Name} & \textbf{Description}\\
\midrule
UC-01 & Propeller WT campaign &
Load N files from a propeller wind tunnel test, each corresponding to a fixed
$(M, N, \beta)$ combination with $\alpha$ as the sweep variable. Derive
$C_T$, $C_P$, $\eta$, $T_C$ with uncertainty propagation. Plot performance
maps. Export.\\
\addlinespace
UC-02 & Flat-table pivot &
Load a single CSV with all datapoints in acquisition order. Inspect column
statistics. Pivot into a structured VarFrame. Apply the same workflow as
UC-01.\\
\addlinespace
UC-03 & CFD disk post-processing &
Load propeller disk data in polar coordinates $(r,\theta)$. Compute azimuthal
averages and disk-integrated in-plane force. Visualize as polar colormap.
Build a Kriging surrogate.\\
\addlinespace
UC-04 & WT balance campaign &
Load power-off balance data with a custom \code{.itceq} file. Perform
repeatability analysis over a \code{repeat} dimension. Compare two tunnel
entries. Generate PDF report.\\
\addlinespace
UC-05 & $C_p$ distribution &
Load surface $C_p$ data indexed by chord-wise position $x/c$, span station,
$\alpha$, $M$. Plot $C_p$ distributions. Compare with wind tunnel pressure
tap data.\\
\addlinespace
UC-06 & Data version comparison &
Load two versions of the same dataset (e.g.\ before and after a tunnel
correction update). Compute delta, delta-\%, and sigma-normalized
differences. Export a comparison report.\\
\addlinespace
UC-07 & Aerospace performance &
Compute takeoff distance, climb gradient, or operational envelope from
aerodynamic and propulsion VarFrame inputs. Output is a VarFrame where every
timestep or flight condition is a datapoint with full computation memory.\\
\addlinespace
UC-08 & EUB: trim deflection band &
Assign wind tunnel balance uncertainties to aerodynamic coefficients,
including correlations from a common calibration, propagate through static
stability analysis, and obtain the Engineering Uncertainty Band on
horizontal tail deflection required for trim.\\
\addlinespace
UC-09 & Multi-fidelity dataset &
Combine WT and CFD data tagged by a fidelity dimension. Optionally fit a
multi-fidelity surrogate to obtain a unified, denser dataset for downstream
performance and stability analyses.\\
\bottomrule
\end{tabularx}
\end{table}
